\documentclass{article}
\usepackage[a4paper, margin=1.5cm]{geometry}
\usepackage{tikz}
\usepackage{pgfplots, pgfplotstable}

\usepackage{amsmath}

\usepackage{siunitx}
\DeclareSIUnit\someunit{unit}

\setlength{\parindent}{0pt}
\setlength{\parskip}{1em}

\title{Least Squares Method - A Quick Overview}
\author{Peleg Sapir}

\begin{document}
\maketitle

\section{Stating the Problem}
Consider the following set of points gathered in some experiment:

\begin{center}
	\begin{tikzpicture}
		\begin{axis}[
				width=12cm,
				axis x line=middle,
				axis y line=middle,
				xlabel=$x\SI{}{[\someunit]}$,
				ylabel=$y\SI{}{[\someunit]}$,
				every axis x label/.style={at={(ticklabel* cs:1.05)}, anchor=west},
				every axis y label/.style={at={(ticklabel* cs:1.05)}, anchor=south},
				axis line style={stealth-stealth, very thick},
				label style={font=\large},
				tick label style={font=\large},
				samples=100,
				xmin=-0.5, xmax=10,
				ymin=-0.9, ymax=20,
				grid=major,
				major grid style={densely dotted, black!15},
			]
			\addplot[only marks, mark=*, mark size=1pt, blue] table [x=x, y=y, col sep=comma] {data.csv};
		\end{axis}
	\end{tikzpicture}
\end{center}

The distribution of the points seems linear, and we would like to know the approximate relation between the dependant variable $y$ and the independant variable $x$.

How can we detemine the \textit{best approximate fit} for all the given points? Here are some examples of possible line fits, with their corresponding equations:

\begin{center}
	\begin{tikzpicture}
		\begin{axis}[
				width=12cm,
				axis x line=middle,
				axis y line=middle,
				xlabel=$x\SI{}{[\someunit]}$,
				ylabel=$y\SI{}{[\someunit]}$,
				every axis x label/.style={at={(ticklabel* cs:1.05)}, anchor=west},
				every axis y label/.style={at={(ticklabel* cs:1.05)}, anchor=south},
				axis line style={stealth-stealth, very thick},
				label style={font=\large},
				tick label style={font=\large},
				samples=100,
				xmin=-0.5, xmax=10,
				ymin=-0.9, ymax=20,
				grid=major,
				major grid style={densely dotted, black!15},
				legend pos=north west,
				legend style={xshift=10pt, nodes={text width=2.7cm, align=left}},
			]
			\addplot[only marks, mark=*, mark size=1pt, blue] table [x=x, y=y, col sep=comma] {data.csv};
			\tikzset{linefit/.style={very thick, dashed, domain={0:10}, opacity=0.5}}
			\addplot[linefit, red] {x+1.5};
			\addplot[linefit, green] {2*x+1};
			\addplot[linefit, purple] {1.5*x+1.2};
			\addplot[linefit, orange] {2.5*x+1.3};
			\addlegendentry{"Real" data}
			\addlegendentry{$a=1.0,\ b=1.5$}
			\addlegendentry{$a=2.0,\ b=1.0$}
			\addlegendentry{$a=1.5,\ b=1.2$}
			\addlegendentry{$a=2.5,\ b=1.3$}
		\end{axis}
	\end{tikzpicture}
\end{center}

\section{Finding the Best Fit}
\subsection{What we're minimizing?}
The line of best fit has the formula $y=mx+b$, and we want to find the coefficients $m$ and $b$. It's impossible to have a single line that would fit all points\footnote{assuming it has zero width, of course.}. Instead, we will try to find the line that minimizes the total $y$-distances between the data points and the line itself.

(HERE BE DRAWING)

% a = 1.66, b = 2.04
% \begin{center}
% 	\begin{tikzpicture}
% 		\begin{axis}[
% 				width=12cm,
% 				axis x line=middle,
% 				axis y line=middle,
% 				xlabel=$x\SI{}{[\someunit]}$,
% 				ylabel=$y\SI{}{[\someunit]}$,
% 				every axis x label/.style={at={(ticklabel* cs:1.05)}, anchor=west},
% 				every axis y label/.style={at={(ticklabel* cs:1.05)}, anchor=south},
% 				axis line style={stealth-stealth, very thick},
% 				label style={font=\large},
% 				tick label style={font=\large},
% 				samples=100,
% 				xmin=-0.5, xmax=10,
% 				ymin=-0.9, ymax=20,
% 				grid=major,
% 				major grid style={densely dotted, black!15},
% 				legend pos=north west,
% 				legend style={xshift=10pt, nodes={text width=2.7cm, align=left}},
% 			]
% 			\pgfplotstableread[col sep=comma]{data.csv}\datatable
%
% 			\def\m{1.66}
% 			\def\b{2.04}
%
% 			\addplot[only marks, mark=*, mark size=1pt, blue] table [x=x, y=y, col sep=comma] {\datatable};
% 			\tikzset{linefit/.style={very thick, dashed, domain={0:10}, opacity=0.5}}
% 			\addplot[linefit, red] {\m*x+\b};
% 			\addlegendentry{"Real" data}
% 			\addlegendentry{Best fit}
% 		\end{axis}
% 	\end{tikzpicture}
% \end{center}

In equation form, each point has the form
\begin{equation}
	p_{i} = \left(x_{i}, y_{i}\right),
	\label{eq:point_form}
\end{equation}
where there are $n$ points (i.e. $i$ goes from $1$ to $n$).

The vertical distance between each point and the line $y=mx+b$ is
\begin{equation}
	r_{i} = y_{i} - \left(mx_{i}+b\right),
	\label{eq:distance_point_line_best_fit}
\end{equation}
and the expression we want to minimize is the sum of the squares of these distances:
\begin{equation}
	S = \sum\limits_{i=1}^{n}r_{i}^{2}.
	\label{eq:sum_to_minimize}
\end{equation}
(we will see later why this expression is a good candidate for minimization)

\subsection{Surprise appearance of linear algebra}
If the line fit \textit{every point}, we could write each of the points as
\begin{equation}
	mx_{i} + b = y_{i}.
	\label{eq:perfect_fit_expression}
\end{equation}

This is a set of $n$ linear equations, which we can write as the following matrix-vector equation:
\begin{equation}
	A\vec{v} = \vec{y},
	\label{eq:}
\end{equation}
or in explicit form:
\begin{equation}
	\begin{bmatrix}
		x_{1}  & 1      \\
		x_{2}  & 1      \\
		x_{3}  & 1      \\
		\vdots & \vdots \\
		x_{n}  & 1      \\
	\end{bmatrix}
	\begin{bmatrix}
		m \\ b
	\end{bmatrix}
	=
	\begin{bmatrix}
		y_{1}  \\
		y_{2}  \\
		y_{3}  \\
		\vdots \\
		y_{n}  \\
	\end{bmatrix}
	\label{eq:matrix_form}
\end{equation}

However, this system has no solution (otherwise we could find the \textit{perfect} fit), and so we will go for the next best thing: finding the projection of $\vec{v}$ onto the subspace spanned by $A$, which we would call $\tilde{v}$. This $\tilde{v}$ would be the minimization of the distance between the $\vec{y}$ and $A\tilde{x}$, i.e. $\min\left\{\left\|\vec{y}-A\tilde{x}\right\|\right\}$, which is the same as finding $\min\left\{\left\|\vec{y}-A\tilde{x}\right\|^{2}\right\}$ (since the square function is an injection over $x\geq0$).

The square norm $\left\|\vec{y}-A\tilde{x}\right\|$ has the explicit form
\begin{equation}
	\left\|\vec{y}-A\tilde{x}\right\|^{2} = \left(y_{1}-u_{1}\right)^{2} + \left(y_{2}-u_{2}\right)^{2} + \dots + \left(y_{n}-u_{n}\right)^{2} = S.
	\label{eq:minimization_as_vector_norm}
\end{equation}

TW: the projection part.
\end{document}
